% Pro T. Annio Milone (pro-milone) --- English translation
% Translated by Claude (Anthropic) from the Latin (Perseus).
% Style guide: STYLE.md.

\ciceroSection{1} Although I fear, gentlemen, that it may be shameful to begin in fear when I speak on behalf of the bravest of men, and most unseemly that --- when T. Annius himself is troubled more for the safety of the commonwealth than for his own --- I should be unable to bring to his cause an equal greatness of spirit, still the strange aspect of this strange court alarms my eyes, which, wherever they fall, look in vain for the old custom of the Forum and the former usage of the courts. For your assembly is not ringed about with a circle of bystanders, as it used to be; we are not hedged in by the accustomed throng;

\ciceroSection{2} and those guards which you see stationed before all the temples, even though they are set there against violence, nonetheless bring some terror to the speaker, so that in the Forum and in a court, although we are fenced about with safeguards salutary and necessary, we still cannot so much as feel no fear without some fear. Were I to think these things ranged against Milo, I would yield to the times, gentlemen; for I would not suppose that, amid such a force of arms, there was any room for a speech. But I am restored and refreshed by the counsel of Cn. Pompeius, a man most wise and most just, who surely would not think it consistent with his own justice to hand over to the weapons of soldiers the same man whom he had handed over to the verdicts of jurors, nor consistent with his wisdom to arm the rashness of an inflamed multitude with public authority.

\ciceroSection{3} Therefore those arms, those centurions, those cohorts proclaim to us not peril but protection, and urge us to be not only of a quiet but even of a great spirit; nor do they promise mere support to my defence, but silence as well. As for the rest of the crowd --- the part of it, at least, that consists of citizens --- it is wholly ours; nor is there any one of those whom you see gazing from every quarter, wherever any part of the Forum can be glimpsed, and awaiting the outcome of this trial, who does not, while he favours the courage of Milo, judge that today there is being fought out a contest over himself, his children, his country, his fortunes. There is one class of men adverse and hostile to us: those whom the madness of P. Clodius fed on plunder and arson and every public ruin; men who were even goaded at yesterday's assembly to dictate to you with their cries what your verdict should be. If perchance there should be any shouting from them, it ought to warn you to keep that citizen who has always, for the sake of your safety, made nothing of that breed of men and of their loudest clamours.

\ciceroSection{4} Therefore be present in spirit, gentlemen, and lay aside whatever fear you have. For if ever you had the power of giving judgement concerning good and brave men, if ever concerning citizens of high desert, if in short there was ever given to chosen men of the most august orders an occasion to declare by deed and verdict the goodwill toward brave and good citizens which they had often signified by look and word, then surely at this present time you hold that power entire: to determine whether we, who have always been devoted to your authority, are always to mourn in misery, or whether, long harried by the most abandoned of citizens, we are at last, through you and through your loyalty, your courage and your wisdom, to be restored.

\ciceroSection{5} For what can be said or imagined more laborious, gentlemen, more anxious, more sorely tried than the case of us two, who, drawn into public life by the hope of the most ample rewards, cannot be free of the dread of the cruellest punishments? For my part, I always thought that the other tempests and storms were to be braved by Milo only in those waves of the popular assemblies, because he had always taken the side of the good against the wicked; but in a court of law, and in that council in which the most august men, drawn from the united orders, sit as judges, I never supposed that Milo's enemies would have any hope at all of destroying not merely his safety but even his glory through the agency of such men as these.

\ciceroSection{6} And yet in this case, gentlemen, we shall not, for the defence against this charge, lean upon the tribunate of T. Annius and all that he has done for the safety of the commonwealth. Unless you see with your own eyes that an ambush was laid for Milo by Clodius, we shall neither beg you to pardon us this charge on account of his many splendid services to the commonwealth, nor demand that, because the death of P. Clodius was your salvation, you should therefore set it down to the credit of Milo's courage rather than to the good fortune of the Roman people. But if his ambush is made clearer than this daylight, then at last I shall implore and entreat you, gentlemen, that, if we have lost everything else, this much at least be left us: that it be lawful to defend our life, with impunity, against the audacity and the weapons of our enemies.

\ciceroSection{7} But before I come to that part of my speech which properly belongs to your inquiry, I think I must refute the things that have so often been flung about, both in the Senate by his enemies and in the assembly by the wicked and, a little while ago, by the prosecutors, so that, every error being cleared away, you may see plainly the matter that comes before your court. They say that it is contrary to divine law for a man to look upon the light who confesses that a man has been slain by him. In what city, then, do men of utter folly maintain this? Why, in the very one which first saw the trial, on a capital charge, of M. Horatius, that bravest of men, who --- though the state was not yet free --- was nonetheless acquitted by the assembly of the Roman people, when he confessed that his sister had been killed by his own hand.

\ciceroSection{8} Or is there anyone who does not know that, when there is an inquiry into a man's being slain, either it is usually denied that the deed was done at all, or it is defended as having been done rightly and lawfully? Unless, that is, you suppose that P. Africanus was out of his mind, who, when he was asked by C. Carbo, tribune of the plebs, seditiously in the assembly, what he thought of the death of Ti. Gracchus, answered that he seemed to have been killed by right. For neither could that famous Servilius Ahala, nor P. Nasica, nor L. Opimius, nor C. Marius, nor, in my own consulship, the Senate be held free of crime, if it were a crime to put wicked citizens to death. And so it is not without reason, gentlemen, that the most learned of men have handed down to memory, even in their invented tales, that the man who, to avenge his father, had slain his mother was acquitted --- when the votes of men were divided --- by the verdict not only of a divinity but of the wisest of goddesses.

\ciceroSection{9} But if the Twelve Tables willed that a nocturnal thief might be killed with impunity by any means whatever, and a thief by day if he should defend himself with a weapon, who is there who would think that, in whatever manner a man is killed, there is ground for punishment --- when he sees that sometimes a sword is held out to us for the slaying of a man by the laws themselves? And yet, if there is any occasion at all for the lawful killing of a man --- and there are many --- surely that one is not merely just but even necessary, when violence is repelled by violence. When a military tribune in the army of C. Marius, a kinsman of that commander, would have ravished the chastity of a soldier, he was killed by the very man on whom he was attempting that violence; for the upright young man chose rather to act at his own peril than to suffer disgrace. And that great man freed him, guiltless of crime, from all peril.

\ciceroSection{10} But against an ambusher and a brigand, what death can be inflicted unjustly? What is the meaning of our escorts, what of our swords? Surely it would not be lawful to have these, if it were on no terms lawful to use them. This, then, gentlemen, is a law not written but born in us --- one which we did not learn, receive, or read, but which we seized, drank in, and wrung from nature herself; a law to which we were not trained but made, not schooled but steeped: that, if our life should fall into any ambush, into the violence and weapons of brigands or of enemies, every honourable means of securing our safety should be lawful.

\ciceroSection{11} For amid arms the laws fall silent, and bid no one wait upon them, since the man who would wait must pay an unjust penalty before he can exact a just one. And yet, with great wisdom and in a certain sense tacitly, the law itself grants the power of self-defence: for it forbids not the killing of a man, but the carrying of a weapon for the purpose of killing a man, so that, when the purpose and not the weapon was at issue, the man who had used a weapon for the sake of defending himself should be judged to have carried the weapon not for the sake of killing a man. Let this, then, stand fast in the case, gentlemen; for I do not doubt that I shall make my defence good to you, if you bear in mind what you cannot forget --- that an ambusher may be killed by right.

\ciceroSection{12} There follows the point most often urged by Milo's enemies: that the Senate judged the bloodshed in which P. Clodius was killed to have been done against the commonwealth. But that bloodshed the Senate approved not by its votes alone but even by its enthusiasm. For how often was that cause pleaded by us in the Senate, and with what assent of the whole order --- assent neither silent nor concealed! When, in a Senate of the fullest attendance, were four men found, or five at most, who did not approve the cause of Milo? This is declared by those half-dead harangues of that singed tribune of the plebs, in which day after day he invidiously charged me with my own power, when he kept saying that the Senate decreed not what it thought but what I wished. If indeed this is to be called power, rather than a modest authority in good causes that comes of great services to the commonwealth, or some little favour among good men that comes of these dutiful labours of mine --- then let it be called so, by all means, provided only that we use it for the safety of the good against the frenzy of the abandoned.

\ciceroSection{13} But this special inquiry, though it is not unjust, the Senate nonetheless never thought ought to be set up; for there were laws, there were standing courts, both for bloodshed and for violence, nor did the death of P. Clodius bring the Senate so great a sorrow and grief that a new inquiry should be established. For when the Senate had been stripped of the power of decreeing an inquiry into that incestuous defilement of his, who can believe that it thought a new inquiry ought to be set up into his death? Why, then, did the Senate decree that the burning of the Senate-house, the assault upon the house of M. Lepidus, this very bloodshed had been done against the commonwealth? Because no violence has ever been undertaken among citizens in a free state that was not against the commonwealth ---

\ciceroSection{14} for no defence against violence is ever to be wished for, but sometimes it is necessary --- unless, that is, the day on which Ti. Gracchus was killed, or the one on which Gaius was, or the arms of Saturninus, even though they were put down for the good of the commonwealth, did not nonetheless wound the commonwealth. And so I myself moved, when it was established that bloodshed had been done on the Appian Way, not that the man who had defended himself had acted against the commonwealth, but, since there was violence and ambush in the matter, I reserved the charge for trial, while I set my mark upon the deed. But if it had been permitted, through that raving tribune of the plebs, for the Senate to carry out what it thought, we should have no new inquiry. For it was decreeing that the matter be tried under the old laws, only out of the ordinary course. The motion was split at the demand of someone or other --- for there is no need for me to bring forward all men's shameful acts --- and so the remaining authority of the Senate was abolished by a bought veto.

\ciceroSection{15} ``But,'' it is said, ``Cn. Pompeius, by his bill, gave judgement both upon the matter and upon the cause: for he brought in a measure concerning the bloodshed done on the Appian Way, in which P. Clodius was killed.'' What, then, did he bring in? Why, that there be an inquiry. And what is to be inquired into? Whether the deed was done? But that is established. By whom? But that is plain. He saw, then, that a defence on the ground of right could be undertaken even in the confession of the deed. And had he not seen that the man who confessed could be acquitted --- seeing that we confess --- he would never have ordered an inquiry at all, nor would he have given you, in your judging, this letter of acquittal, so salutary, any more than that grim one of condemnation. To my mind Cn. Pompeius has not only judged nothing more severely against Milo, but has even laid down what you, in judging, ought to look to. For the man who appointed to confession not a penalty but a defence thought that the cause of the death, and not the death itself, was to be inquired into.

\ciceroSection{16} As for whether he thought what he did of his own accord was to be set down to Publius Clodius or to the times --- that surely he himself will say. M. Drusus, tribune of the plebs, a most noble man, in his own day the champion of the Senate and, in those times indeed, almost its patron, the maternal uncle of this juror of ours, the bravest of men, M. Cato, was killed in his own house. Nothing was laid before the people about his death; no inquiry was decreed by the Senate. What grief there was in this city, as we have heard from our fathers, when that nocturnal violence was brought against P. Africanus as he rested in his own house! Who did not groan then, who did not burn with anguish that the death of a man whom all would have wished, if it could be, to be immortal, was not so much as awaited as a natural end? Was there, then, any inquiry brought concerning the death of Africanus? Surely none.

\ciceroSection{17} Why so? Because famous men are not killed by one crime and obscure men by another. Let there be a difference between the dignity of the lives of the highest and the lowest; but death inflicted by crime, at least, let it be held bound by the same penalties and the same laws. Unless, perhaps, a man will be the more a parricide if he kills a father of consular rank than if he kills a lowly one, or the death of P. Clodius will be the more atrocious for this, that he was killed among the monuments of his own ancestors --- for so it is often said by those men --- just as though Appius the Blind built that road not for the people to use, but for his own descendants to play the brigand upon with impunity!

\ciceroSection{18} And so, when on that same Appian Way P. Clodius had killed M. Papirius, a most distinguished Roman knight, that deed was not to be punished --- for a noble man had killed a Roman knight among his own monuments --- but now what tragedies the name of that same Appian Way calls up! The road that was passed over in silence when it was stained earlier with the bloodshed of an honourable and innocent man is the very road now invoked again and again, since it has been steeped in the blood of a brigand and a parricide. But why do I recall those things? A slave of P. Clodius was caught in the temple of Castor, whom he had stationed there to kill Cn. Pompeius. A dagger was wrenched from his hands as he confessed. Thereafter Pompeius kept away from the Forum, kept away from the Senate, kept away from public life; he sheltered himself behind his door and his walls, not behind the right of laws and courts.

\ciceroSection{19} Was any bill brought, was any new inquiry decreed? And yet, if ever matter, or man, or time was worthy of one, surely all these things in that case were of the highest. An ambusher had been stationed in the Forum and in the very forecourt of the Senate; death, moreover, was being prepared for the man on whose life the safety of the state rested; and that, too, at such a season of the commonwealth that, had that one man perished, not this state alone but all nations would have fallen. Unless, indeed, because the deed was not accomplished, it was therefore not to be punished --- just as though it were the outcomes of things, and not the designs of men, that the laws avenge. There was less to grieve over, the deed not being accomplished, but it was no less surely to be punished.

\ciceroSection{20} How often have I myself, gentlemen, escaped the weapons of P. Clodius and his bloody hands! Had not either my own fortune or the commonwealth's preserved me from them, who in the end would have brought an inquiry concerning my death? But we are fools to dare to set Drusus, Africanus, Pompeius, our very selves, alongside P. Clodius. Those losses were bearable: the death of P. Clodius no one can bear with an even mind. The Senate mourns, the equestrian order grieves, the whole state is worn out with age, the towns are in squalor, the colonies are stricken, the very fields, in short, miss so beneficent, so salutary, so gentle a citizen.

\ciceroSection{21} That was not the reason, gentlemen, surely it was not, why Pompeius should think an inquiry ought to be brought; but a wise man, endowed with a certain lofty and divine mind, saw many things: that Clodius had been his enemy and Milo his friend; that if he too should rejoice in the common gladness of all, the loyalty of a reconciled goodwill might seem the more shaky. He saw many other things besides, but this above all: that, however harshly he himself had taken it, you would nonetheless judge bravely. And so he chose, from the most flourishing orders, their very luminaries; nor indeed, as certain men keep saying, did he set aside my friends in the choosing of the jurors. For neither did that most just of men plan such a thing, nor, in choosing good men, could he have achieved it, even had he wished. For my influence is not bounded by close intimacies, which cannot spread wide, since the daily commerce of life cannot be carried on with many; but if I have any influence, I have it from this, that the commonwealth has joined me with good men. And when he was choosing from among these the best men, and judged that this above all bore upon his own good faith, he could not but choose men devoted to me.

\ciceroSection{22} As for his wishing above all that you, L. Domitius, should preside over this inquiry, he sought nothing else but justice, gravity, humanity, good faith. He brought it in that the presiding officer must be of consular rank: because, I take it, he held it the part of leading men to withstand both the fickleness of the multitude and the rashness of the abandoned. From among the consulars he chose you in preference to all: for you had given, ever since your youth, the greatest proofs of how thoroughly you despise the frenzies of the mob.

\ciceroSection{23} Wherefore, gentlemen --- that we may come at last to the case and the charge --- if neither is every confession of a deed without precedent, nor has anything been judged by the Senate concerning our cause otherwise than we ourselves should wish, and the very proposer of the law, though there was no dispute about the fact, nonetheless willed that there be a debate about the right, and jurors of such a kind have been chosen, and a man set over the inquiry who will weigh these matters justly and wisely --- it remains, gentlemen, that you ought now to inquire into nothing else but which of the two laid an ambush for the other. That you may the more easily perceive this by the proofs, while I lay the matter before you briefly, I beg you to attend with care.

\ciceroSection{24} P. Clodius, when he had resolved to harry the commonwealth with every crime during his praetorship, and saw that the elections of the previous year had been so dragged out that he could not hold the praetorship for many months --- since he had no eye to the rank of the office, as others have, but both wished to escape L. Paulus as his colleague, a citizen of singular virtue, and sought a whole year for tearing the commonwealth to pieces --- suddenly abandoned his own year and transferred himself to the next, not, as happens, out of any religious scruple, but in order to have, as he himself put it, for the conduct of his praetorship --- that is, for the overthrow of the commonwealth --- a full and entire year.

\ciceroSection{25} It occurred to him that his praetorship would be crippled and feeble with Milo as consul; and, what is more, he saw that Milo was being made consul by the highest agreement of the Roman people. He betook himself to Milo's competitors, but in such a way that he alone steered the whole canvass, even against their will, that he carried the whole election, as he kept boasting, on his own shoulders. He summoned the tribes, thrust himself in, enrolled a new Colline tribe by a levy of the most abandoned citizens. The more confusion that man stirred up, the more this man grew stronger by the day. When that fellow, ready for every crime, saw the bravest of men, his bitterest enemy, the most certain of consuls, and understood that this had been declared again and again not only by the talk but even by the votes of the Roman people, he began to act openly and to say plainly that Milo must be killed.

\ciceroSection{26} He had brought down from the Apennines the rustic and barbarous slaves with whom he had ravaged the public forests and harried Etruria --- the very men you have been seeing. The matter was not in the least obscure. For he kept saying openly that the consulship could not be wrested from Milo, but his life could. He hinted at this often in the Senate, said it in the assembly; nay, even to M. Favonius, the bravest of men, who asked him on what hope he was raging while Milo lived, he answered that within three days, or four at most, Milo would be dead --- a saying of his which Favonius at once reported to this M. Cato.

\ciceroSection{27} Meanwhile, since Clodius knew --- and this was no hard thing to learn from the people of Lanuvium --- that on the thirteenth day before the Kalends of February Milo had a journey to make to Lanuvium, a journey set, prescribed, and unavoidable, to nominate the flamen, because Milo was dictator of Lanuvium, Clodius himself set out from Rome suddenly, the day before, in order --- as the event made plain --- to lay an ambush for Milo before his own estate; and he set out in such a way as to leave behind a turbulent public meeting, held on that very day, at which his frenzy was sorely missed --- a meeting he would never have left, had he not wished to keep the appointed place and hour of his crime.

\ciceroSection{28} Milo, on the other hand, having been in the Senate that day until the Senate was dismissed, came home, changed his shoes and his clothes, waited a little while, as one does, for his wife to make herself ready, and then set out at the hour by which Clodius, had he indeed meant to come to Rome that day, could already have returned. Clodius meets him --- lightly equipped, on horseback, no carriage, no baggage, no Greek companions in his usual way, without his wife, which was almost unheard of --- while this so-called ambusher, who had supposedly arranged that journey for the purpose of committing murder, was riding in a carriage with his wife, wrapped in a travelling-cloak, with a great, encumbering, womanish, and dainty train of maids and boys.

\ciceroSection{29} He falls in with Clodius before Clodius's estate, at about the eleventh hour, or not much off it. At once a number of men with weapons rush upon him from the higher ground; head-on, they kill the driver. But when Milo, throwing back his cloak, had leapt down from the carriage and was defending himself with a fierce spirit, the men who were with Clodius drew their swords and, some of them running back to the carriage to attack Milo from behind, others --- because they supposed Milo already killed --- beginning to cut down those of his slaves who were behind; of these, those who were loyal and steadfast in spirit toward their master, some were slain, while others, when they saw the fighting at the carriage, were kept from coming to their master's aid, heard from Clodius himself that Milo was killed, and truly believed it, did a thing --- I shall say it openly, not to divert the charge, but because it is what happened --- that Milo's slaves did, with their master neither ordering nor knowing nor present, the very thing that every man would have wished his own slaves to do in such a case.

\ciceroSection{30} These things were done just as I have set them out, gentlemen: the ambusher was overcome, force was conquered by force, or rather audacity was crushed by valour. I say nothing of what the commonwealth gained, nothing of what you gained, nothing of what all good men gained: let none of that profit Milo, who was born under this destiny, that he could not even save himself without at the same time saving the commonwealth and you. If this could not be done lawfully, I have nothing to plead. But if reason has prescribed this to the learned, necessity to the barbarian, custom to the nations, and nature herself to the very beasts --- that they should always repel all violence, by whatever means they can, from their body, their head, their life --- then you cannot judge this deed wicked without judging at the same time that all who fall among brigands must perish, either by those men's weapons or by your votes.

\ciceroSection{31} But if Milo had thought so, surely it would have been more to be wished, for him, to offer his throat to P. Clodius --- a throat sought by Clodius more than once, and not then for the first time --- than to be slaughtered by you because he had not handed himself over to Clodius for slaughter. But if none of you thinks so, then this is what now comes to trial: not whether Clodius was killed, which we admit, but whether rightly or wrongly, a question often raised in many cases. That an ambush was laid is agreed, and it is this that the Senate judged to have been committed against the commonwealth; by which of the two it was laid is uncertain. It is on this point, then, that the bill was carried, that an inquiry be held. So the Senate censured the act, not the man, and Pompey raised the question of the law, not of the fact. Does anything else, then, come to trial except which of the two laid the ambush for the other? Surely nothing: if Milo for Clodius, let him not go unpunished; if Clodius for Milo, then let us be acquitted of the crime.

\ciceroSection{32} By what means, then, can it be proved that Clodius laid the ambush for Milo? With a beast so bold, so wicked, it is enough to show that he had great cause, great hope set before him in Milo's death, great advantages to gain. And so let that maxim of Cassius hold good in these two persons --- to whose benefit a thing was done --- even though good men are driven into crime by no profit, and the wicked often by a small one. Now, with Milo dead, Clodius would attain this: not only that he should be praetor under no consul able to keep him from doing any wrong, but also that he should be praetor under consuls who, if not helping him, would at any rate connive, so that he might hope to play his part in those frenzies he had planned --- whose attempts, as he reckoned himself, those consuls would not wish to suppress, even if they could, supposing themselves to owe him so great a service, and, even if they wished to, could perhaps scarcely break the audacity of a man most criminal, hardened now by long habit.

\ciceroSection{33} Or are you alone in your ignorance, gentlemen? Do you dwell in this city as strangers? Do your ears go travelling abroad, and have no part in this talk that runs through the whole community --- what laws, if laws they are to be called and not firebrands of the city, plagues of the commonwealth, that man meant to fasten and brand upon us all? Produce, produce, I beg you, Sextus Clodius, that little book of your laws, which they say you snatched from his house and carried off out of the very midst of arms and the nightly mob, like the Palladium, so that you might bring this splendid gift, this equipment of a tribunate, to someone --- had you found him --- to administer the tribunate at your direction. And he gave me a look, those very eyes he used to wear when he was threatening everyone with everything. That light of the Senate-house moves me, to be sure! What? Do you think me angry with you, Sextus, when you have punished my bitterest enemy far more cruelly than the kindness of my nature would have demanded? You cast P. Clodius's bloody corpse out of the house; you flung it into the street; you stripped it of its ancestral images, its funeral rites, its procession, its eulogy, and left it half-burned on a wretched heap of sticks, to be torn by the dogs of the night. Therefore, though you acted abominably, still, because it was upon my enemy that you spent your cruelty, I cannot praise you, but I certainly ought not to be angry.

\ciceroSection{34} You have heard, gentlemen, how much it was in Clodius's interest that Milo be killed: now turn your minds in turn to Milo. What interest had Milo in Clodius being killed? What reason was there for Milo --- I will not say to allow it, but to wish for it? Clodius stood in Milo's way as he hoped for the consulship. But with Clodius opposing, Milo's hope went forward all the same --- indeed it went forward the more for that, and he had no better canvasser than Clodius. With you, gentlemen, the memory of Milo's services toward me and toward the commonwealth carried weight; my prayers and tears carried weight, by which I felt you then to be wonderfully moved; but far more weight was carried by the fear of impending dangers. For what citizen was there who could face P. Clodius's praetorship, set loose, without the greatest dread of revolution? And you saw that it would be set loose, unless there were a consul with the daring and the power to bind it fast. When the whole Roman people felt that Milo alone was that man, who would hesitate to free himself from fear, and the commonwealth from peril, by his vote? But now, with Clodius removed, Milo must strive by the ordinary means to guard his standing; that singular glory, granted to him alone, which grew greater each day by the breaking of Clodius's frenzies, has now fallen with Clodius's death. You have won the right to fear no citizen; he has lost the exercise of his valour, his support for the consulship, the perennial spring of his glory. And so Milo's consulship, which could not be shaken while Clodius lived, has only now, with Clodius dead, begun to be assailed. Clodius's death, then, not only does Milo no good, but actually does him harm.

\ciceroSection{35} But hatred prevailed; he acted in anger, he acted as an enemy; he was the avenger of his own wrong, the punisher of his own grievance. What? If these motives were --- I will not say greater in Clodius than in Milo, but greatest in Clodius, none at all in Milo, what more do you want? For why should Milo have hated Clodius, the crop and the material of his glory, beyond that civic hatred with which we hate all wicked men? In Clodius there was reason to hate Milo: first, the defender of my safety; then the harrier of his frenzy, the tamer of his arms; last of all, his very prosecutor; for Clodius was, as long as he lived, a defendant of Milo's under the Lex Plotia (a law on public violence). With what spirit, then, do you believe that tyrant bore this? How great his hatred, and, in a man so unjust, how even just it was?

\ciceroSection{36} It remains that his own nature and habit should now defend Clodius, while these same things convict Milo. Clodius never did anything by violence; Milo, everything by violence. What? When I, gentlemen, while you mourned, withdrew from the city, was it a trial I feared --- and not slaves, not arms, not violence? What just cause, then, would there have been for restoring me, had there not been an unjust one for casting me out? He had named a day for me, I suppose, had proposed a fine, had brought an action of high treason; and no doubt I had a trial to fear, in a case either bad, or merely my own, and not the most glorious and yours as well. No --- I was unwilling that my fellow citizens, men saved by my counsels and my dangers, should be thrown into peril for my sake, against the arms of slaves, of beggared citizens, and of criminals.

\ciceroSection{37} For I saw --- yes, I saw with my own eyes --- this very Q. Hortensius, the light and ornament of the commonwealth, almost killed by a band of slaves, while he was standing by me; and in that throng C. Vibienus, a senator, a man of the highest worth, who was there together with him, was so mauled that he lost his life. And so, when did that dagger of his, the one he had received from Catiline, ever afterward rest? It was aimed at me; I would not let you be exposed to it for my sake; it lay in wait for Pompey; it stained the Appian Way, the monument of his own name, with the murder of Papirius; this same dagger, after a long interval, was turned again upon me; lately indeed, as you know, it nearly finished me near the Regia.

\ciceroSection{38} What was there like this in Milo? His whole violence was always to this end: that P. Clodius, since he could not be dragged into court, should not hold the state in his grip, crushed by force. Had Milo wished to kill him, what great occasions there were, and how often, and how splendid! Could he not have avenged himself lawfully, when he was defending his house and his household gods with Clodius storming them? Could he not, when an excellent citizen and most brave man, his colleague P. Sestius, was wounded? Could he not, when Q. Fabricius, a man of the highest worth, was driven off as he was carrying a law for my recall, and a most cruel slaughter was done in the Forum? Could he not, when the house of L. Caecilius, that most just and brave praetor, was stormed? Could he not, on that day when the law concerning me was passed, when the gathering of all Italy, which my safety had stirred up, would gladly have acknowledged the glory of that deed --- so that, even if Milo had done it, the whole community would have claimed the praise as its own?

\ciceroSection{39} But what a time it was! There was that most illustrious and brave man, the consul, an enemy of Clodius, P. Lentulus, the avenger of that man's crime, the champion of the Senate, the defender of your will, the patron of the public consensus, the restorer of my safety; seven praetors, eight tribunes of the plebs, his adversaries and my defenders; Cn. Pompey, the author and leader of my recall, that man's enemy, whose opinion --- most weighty and most distinguished --- the whole Senate followed concerning my safety, who exhorted the Roman people, who, when he had passed a decree about me at Capua, himself gave the signal to all Italy, eager and imploring his protection, to flock to Rome for my restoration; in short, the hatred of all the citizens blazed against him out of longing for me, so that whoever had then made an end of him would have been thought of not in terms of impunity, but of rewards.

\ciceroSection{40} Yet then Milo held himself in check, and summoned P. Clodius to trial twice, to violence never. What? When Milo was a private citizen and a defendant, with P. Clodius prosecuting him before the people, when an attack was made on Cn. Pompey as he was speaking for Milo, what an opportunity there was then --- not only an opportunity but a reason --- for crushing that man! Lately, indeed, when M. Antonius had brought the highest hope of safety to all good men, and that most noble young man had most bravely taken up the gravest part of the commonwealth's burden, and held that beast already snared, as it tried to slip the nooses of the court --- what an opening, what a moment that was, immortal gods! When the man hid himself, fleeing, in the darkness of a staircase, what a great thing it would have been for Milo to make an end of that pestilence with no odium to himself, and to the very greatest glory of M. Antonius!

\ciceroSection{41} What? At the elections on the Campus, how many times was the chance there! When that man had burst into the voting-enclosures, had arranged for swords to be drawn and stones to be thrown, and then, suddenly terrified at the look on Milo's face, fled toward the Tiber, while you and all good men were offering up prayers that Milo might be pleased to use his valour. The man, then, whom Milo would not kill with everyone's gratitude --- did he wish to kill him to the complaint of some? The man whom he did not dare kill rightly, in the right place, at the right time, with impunity --- did he not hesitate to kill wrongly, in an unfavourable place, at the wrong time, at the risk of his own life?

\ciceroSection{42} Especially, gentlemen, when a contest for the highest office and the day of the elections were close at hand, at which time --- for I know how timid ambition is, and how great and how anxious is the desire for the consulship --- we fear everything, not only the things that can be openly censured but even those that can be secretly imagined; we shudder at a rumour, a tale false, fabricated, trivial; we watch the faces and the eyes of all. For there is nothing so soft, so tender, so fragile or so pliable as the goodwill and the feeling of the citizens toward us, who are not only angry at the misconduct of candidates, but often grow fastidious even at deeds rightly done.

\ciceroSection{43} Did Milo, then, setting before himself this hoped-for and longed-for day of the Campus, come to those august auspices of the centuries with bloody hands, parading and confessing his crime and his outrage? How incredible this is in him --- how, in Clodius, the same thing is not to be doubted, when he supposed that, with Milo killed, he himself would reign! And what --- for here is the very head of audacity, gentlemen --- who does not know that the greatest enticement to wrongdoing is the hope of impunity? In which of the two, then, was this? In Milo, who even now is a defendant for a deed either glorious or at any rate necessary, or in Clodius, who had so despised the courts and their penalties that nothing pleased him which was either permitted by nature's law or allowed by the laws of men?

\ciceroSection{44} But why do I argue, why do I dispute further? I appeal to you, Q. Petilius, the best and bravest of citizens; I call you to witness, M. Cato, you whom a kind of divine chance has given me as jurors. You heard from M. Favonius that Clodius had said to him --- and you heard it while Clodius was alive --- that Milo would die within three days. Three days after he had said it, the thing was done. When that man did not hesitate to disclose what he was planning, can you hesitate as to what he did?

\ciceroSection{45} How, then, did the day not catch Clodius unawares? I have just said. There was no trouble in knowing the fixed sacrifices of the dictator of Lanuvium. He saw that Milo must of necessity set out for Lanuvium on the very day he did set out: and so he forestalled him. But on what day? The day on which, as I said before, there was a most frenzied public meeting, stirred up by his own hireling tribune of the plebs --- a day, a meeting, a clamour that he would never have left, had he not been hastening to a premeditated crime. So for him there was not even a reason for the journey --- there was, indeed, a reason to stay; for Milo there was no possibility of staying, and for setting out not only a reason but a necessity. What if, just as that man knew Milo would be on the road that day, so Milo could not even suspect the same of Clodius?

\ciceroSection{46} First I ask how he could have known it --- a question you cannot put in the same way about Clodius. For even if Clodius had asked no one but T. Patina, his closest friend, he could have learned that on that very day the flamen had to be nominated by Milo, the dictator, at Lanuvium. But there were a great many others from whom he could have learned it most easily: all the people of Lanuvium, of course. From whom did Milo inquire about Clodius's return? Grant that he did inquire --- see what I concede to you --- grant that he even bribed a slave, as my friend Q. Arrius said. Read the testimony of your own witnesses. C. Causinius Schola of Interamna, Clodius's closest friend and companion, testified that P. Clodius meant to stay that day at his Alban villa, but that word came to him suddenly that the architect Cyrus had died, and so on the instant he determined to set out for Rome. C. Clodius, likewise a companion of P. Clodius, said the same.

\ciceroSection{47} See, gentlemen, how much these testimonies have settled. First, Milo is certainly cleared of having set out with the intention of lying in wait for Clodius on the road: for, on this showing, Clodius was not going to fall in with him at all. Next --- for I see no reason not to plead my own business too --- you know, gentlemen, that there were men who, in urging this very bill, said that the murder was done by Milo's hand, but by the design of someone greater. Plainly it was me whom those abject and ruined men were portraying as a brigand and an assassin. They lie prostrate, by their own witnesses, who deny that Clodius would have returned to Rome that day had he not heard about Cyrus. I breathed again; I was set free; I do not fear that I should seem to have planned a thing which I could not even have suspected.

\ciceroSection{48} Now I will pursue the rest; for this occurs to me: so then not even Clodius was thinking of an ambush, since he meant to stay at his Alban villa. Yes --- if he had not been going to leave the villa for murder. For I see that the man who is said to have brought word of Cyrus's death brought no such word, but rather that Milo was drawing near. For what news could he bring of Cyrus, whom Clodius, as he was setting out from Rome, had left dying? I sealed his will along with him, I was present; but he had made his will openly, and had named both Clodius his heir and me. The man whom he had left, the day before, at the third hour, gasping out his life --- was it only the next day, at the tenth hour, that he was finally being told of his death?

\ciceroSection{49} Come, suppose it was so: what reason was there for him to hasten to Rome, to throw himself into the night? What was there in his being heir that called for haste? First, there was no reason why haste was needed; and then, if there were, what, pray, was there that he could gain that night and lose if he came to Rome the next morning? And just as a nightly arrival at the city was for him to be avoided rather than sought, so for Milo, since he was the ambusher --- if he knew that Clodius would approach the city by night --- there was reason to lie in wait and wait for him.

\ciceroSection{50} He would have killed him by night: he would have killed him in a treacherous place full of brigands. No one would have disbelieved him as he denied it, when everyone wishes him safe even as he confesses. The very place would first have borne the charge --- that haunt and harbour of brigands; then the dumb solitude would not have informed against Milo, nor the blind night exposed him; next, many men whom Clodius had abused, robbed, driven from their goods, and many who feared the same, would have fallen under suspicion; in short, all Etruria would have been summoned as defendant.

\ciceroSection{51} And on that very day, Clodius, returning from Aricia, certainly turned aside to his place at the Alban villa. Even if Milo knew that he had been at Aricia, still he ought to have suspected that, even if Clodius wished to return to Rome that day, he would turn aside to his own villa, which bordered the road. Why did he not meet him beforehand, to keep that man from settling in at the villa, nor lie in wait in the place to which Clodius would come by night?

\ciceroSection{52} I see, gentlemen, that thus far everything is consistent: that it was even to Milo's advantage for Clodius to live, while to Clodius, for the things he had set his heart on, the death of Milo was most to be wished; that Clodius's hatred toward Milo was most bitter, Milo's toward Clodius none at all; that Clodius's habit was a perpetual one of offering violence, Milo's only of repelling it; that death was threatened against Milo by Clodius and proclaimed openly, while nothing of the kind was ever heard from Milo; that the day of Milo's setting-out was known to Clodius, while Clodius's return was unknown to Milo; that Milo's journey was a necessary one, Clodius's even rather uncalled-for; that Milo had openly announced he would leave that day, while Clodius had concealed that he would return that day; that Milo had changed his plan in nothing, while Clodius had invented a reason for changing his; that Milo, if he were lying in wait, had to wait through the night near the city, while Clodius, even if he did not fear Milo, still had to dread approaching the city by night.

\ciceroSection{53} Let us now look at what is the head of the matter: which of the two, in the end, that very place where they met was better suited for laying an ambush. Is this, gentlemen, even to be doubted, or pondered any longer? Before Clodius's estate --- an estate where, because of those mad foundation-works, a good thousand able-bodied men were easily kept busy --- did Milo think that from his adversary's raised and lofty ground he would have the upper hand, and for that reason choose that place above all for a fight? Or was he rather waited for in that place by the man who, on the strength of the ground itself, had planned to make the attack? The thing speaks for itself, gentlemen, and that always carries the greatest weight.

\ciceroSection{54} If you were not hearing these things told, gentlemen, but seeing them painted, even so it would be plain which of the two was the plotter, which had no harm in mind --- when the one was riding in a carriage, wrapped in a travelling-cloak, his wife seated at his side. Which of these was not the very picture of encumbrance: the dress, or the vehicle, or the companion? What could be less ready for a fight than a man tangled in his cloak, hampered by his carriage, all but pinned down by his wife? --- Now look at the other: first, leaving his country house, suddenly. Why? In the evening. What need of that? Slowly. How does that fit, especially at that hour? He turns aside to Pompey's villa. To see Pompey? He knew he was at Alsium. To inspect the villa? He had been in it a thousand times. What, then, was it? Delay and stalling: he would not quit the spot until Milo came up.

\ciceroSection{55} Come now, compare the journey of the unencumbered bandit with Milo's baggage. Always before, the one travelled with his wife; this time, without her. Never except in a carriage; this time, on horseback. Greek hangers-on for company wherever he went, even when he was hurrying to his Etruscan camp; this time, no triflers in his retinue. Milo, who never did so, was on this occasion taking along, as it chanced, his wife's musician-boys and a flock of maidservants; the other, who always led with him whores, always catamites, always harlots, this time had no one --- unless you would say it was man picked out, man by man. Why, then, was he beaten? Because the traveller is not always killed by the bandit; now and then the bandit too is killed by the traveller --- because, though Clodius came prepared against men unprepared, Clodius himself was a woman who had fallen among men.

\ciceroSection{56} Nor indeed was Milo ever so unprepared against him as not to be very nearly prepared enough. He always reckoned both how much it mattered to P. Clodius that he himself should perish, and how deeply that man hated him, and how much that man dared. And so, knowing his own life set up for the highest prizes and all but knocked down at auction, he never threw it into danger without an escort, without a guard. Add the chances; add the uncertain outcomes of fights and the impartiality of Mars, who often, when a man is already stripping and exulting, overturns him and strikes him down by the very man he had cast off; add the blundering of a commander dined, drunk, and yawning, who, when he had left an enemy cut off at his rear, gave no thought to that enemy's last companions; and falling among them, kindled with rage and despairing of their master's life, he stuck fast in the punishment which faithful slaves exacted from him for their master's life.

\ciceroSection{57} Why, then, did Milo set them free? Because, no doubt, he feared they might inform, might be unable to bear the pain, might be forced by torture to confess that P. Clodius had been killed by Milo's slaves on the Appian Way. What need of the rack? What are you asking? Whether he killed him? He killed him. Lawfully, or unlawfully? That is nothing to the torturer: for on the rack the inquiry is into the deed, in court into the law. What, then, must be inquired into in the case, let us pursue here; what you wish to find out by torture, that we confess. But if you ask why he set them free, rather than why he rewarded them too little handsomely, you do not know how to fault a deed done by an enemy.

\ciceroSection{58} For this same man who has always said everything firmly and bravely, M. Cato, said --- and said it in a turbulent public meeting, which was nonetheless calmed by his authority --- that those who had defended their master's life were most worthy not only of freedom but of every reward. For what reward is great enough for slaves so kindly, so good, so faithful, through whom he lives? And yet even that is not worth so much as this: that through these same men he did not glut the mind and eyes of a most cruel enemy with his own blood and wounds. Had he not freed them, the preservers of their master, the avengers of crime, the defenders against murder would have had to be handed over even to the torture. But in the midst of these troubles he has nothing that he bears with less distress than this: that, even if anything should happen to himself, the deserved reward has nonetheless been paid them in full.

\ciceroSection{59} But the examinations press hard upon Milo --- those held just now in the Hall of Liberty. Of whose slaves? You ask? Of P. Clodius's. Who demanded them? Appius. Who produced them? Appius. From where? From Appius. Good gods! What could be done more harshly? Clodius came nearer to the gods than when he had broken in upon their very presence, since his death is now inquired into as though sacred rites had been violated. And yet our forefathers refused to allow examination against a master, not because the truth could not be found, but because it seemed shameful, and sadder to the masters than death itself. When examination is made of a prosecutor's slave against the defendant, can the truth be found?

\ciceroSection{60} Come now, what was the examination, and of what sort? ``Here, you, Rufio'' --- for instance --- ``mind you do not lie: did Clodius lay a plot against Milo?'' ``He did'': the cross is certain. ``He laid none'': freedom is hoped for. What is surer than such an examination? Men snatched up suddenly for examination are nonetheless kept apart, the rest, and thrown into cells, that no one may speak with them: but these, after they had been a hundred days in the keeping of the prosecutor, were produced by that very prosecutor. What can be said more honest than such an examination, what more uncorrupted?

\ciceroSection{61} But if you do not yet see clearly enough, when the matter itself shines out with so many and so clear arguments and signs --- that Milo returned to Rome with a mind pure and whole, stained by no crime, terrified by no fear, undone by no guilty conscience --- then recall, by the immortal gods, what was the speed of his return, what his entry into the Forum while the Senate-house was burning, what greatness of spirit, what countenance, what speech. Nor indeed did he commit himself to the people alone, but also to the Senate; nor to the Senate only, but also to the public guards and arms; nor to these only, but to the power of the man to whom the Senate had committed the whole commonwealth, all the manhood of Italy, all the arms of the Roman people. To this man he surely would never have surrendered himself, had he not trusted his own cause --- and that to one who was hearing everything, fearing much, suspecting many things, believing some. Great is the force of conscience, gentlemen, and great in both directions: so that those who have committed no wrong feel no fear, while those who have sinned think the penalty forever set before their eyes.

\ciceroSection{62} Nor indeed without sure reason was Milo's cause always approved by the Senate; the wisest of men saw the reasoning of the deed, his presence of mind, the steadfastness of his defence. Or have you forgotten, gentlemen, when the news of Clodius's death was fresh, the talk and the opinions not only of Milo's enemies but even of some who knew no better? They kept saying he would not return to Rome.

\ciceroSection{63} For if he had done it in an angry and excited spirit --- so as to butcher his enemy in the heat of hatred --- they supposed he had reckoned the death of P. Clodius worth so much that he would bear exile from his country with an even mind, now that he had filled his hatred with the blood of his enemy; or if he had wished by that man's death to free his country, they thought a brave man would not hesitate, having brought safety to the Roman people at his own peril, to yield with an even mind to the laws, to carry off with him everlasting glory, and to leave for your enjoyment what he himself had preserved. Many even talked of Catiline and those monstrous portents: he will break out, he will seize some position, he will make war upon his country. Wretched at times are the citizens who have deserved supremely well of the commonwealth, in whom men not only forget the most splendid deeds but even suspect wicked ones!

\ciceroSection{64} Those things, then, were false --- things that would certainly have proved true, had Milo committed anything he could not honourably and truthfully defend. What of the charges afterward heaped upon him, which would have crushed the conscience even of moderate offences --- how he bore them, immortal gods! Bore them? No, rather how he despised them and counted them as nothing, things that neither a guilty man, however great his spirit, nor an innocent one, unless a most brave man, could have disregarded! Word was given out that a great quantity of shields, swords, javelins, even of bridles, could be seized; they said there was no quarter in the city, no back-alley, in which a house had not been hired for Milo; that arms had been carried down the Tiber to his villa at Ocriculum; that his house on the Capitoline slope was stuffed with shields, all of it full of firebrands made ready to burn the city: these things were not only reported, but almost believed, nor were they rejected until they were investigated.

\ciceroSection{65} I praised, for my part, the incredible diligence of Cn. Pompey; but I will say, gentlemen, what I think. They are forced to hear too many things, and can do no otherwise, those to whom the whole commonwealth has been committed. Why, even a certain Licinius had to be heard, some sacrificer or other from the Circus Maximus, who said that Milo's slaves, having got drunk at his house, confessed to him that they had conspired to murder Cn. Pompey, and that afterward he himself had been run through with a sword by one of them, lest he inform. Word is taken to Pompey in his gardens; I am summoned among the first; on the advice of his friends he lays the matter before the Senate. I could not but be undone by fear, with so grave a suspicion hanging over that guardian of mine and of my country; yet I marvelled all the same that a sacrificer should be believed, the confession of slaves heard, a wound in the side --- which looked like a pin-prick --- accepted as the thrust of a gladiator.

\ciceroSection{66} But, as I understand it, Pompey was guarding against danger rather than fearing it --- not only against the things that were to be feared, but against everything, lest you should fear anything. The house of C. Caesar, that most illustrious and most brave man, was reported to have been besieged for many hours of the night: no one had heard of it in so frequented a place, no one had felt it; yet it was reported. I could not suspect Cn. Pompey, a man of surpassing virtue, of being timorous; no diligence undertaken for the whole commonwealth did I think excessive. In a most crowded Senate, lately on the Capitol, a senator was found to say that Milo was carrying a weapon. He stripped himself in that most sacred temple, since the life of such a citizen and such a man could win no credence, so that, with himself silent, the matter itself might speak. All these things were found to be falsely and spitefully invented; and yet even now Milo is feared.

\ciceroSection{67} It is no longer this Clodian charge that we fear, but your suspicions, Cn. Pompey --- you I address, and in such a voice that you may be able to hear me --- your suspicions, I say, we shudder at. If you fear Milo, if you think that this man either now plots wickedly against your life or has at some time contrived something, if the levy throughout Italy --- as some of your recruiting officers kept repeating --- if these arms, if the Capitoline cohorts, if the sentries, the watches, the chosen young men who guard your person and your house, are armed against an onslaught of Milo's, and all of that has been arrayed, made ready, and aimed against this one man --- then surely great force is judged to be in him, an incredible spirit, the strength and resources of no single man, if indeed against this one man both a most outstanding general has been chosen and the whole commonwealth has been put under arms.

\ciceroSection{68} But who does not understand that all the sick and tottering parts of the commonwealth were committed to you, that with these arms you might heal and steady them? And if Milo had been given the chance, he would surely have proved to you yourself that no man was ever dearer to another than you to him; that he had never shunned any danger for your standing; that against that foulest plague itself he had very often striven for your glory; that his tribunate had been steered by your counsels toward my safety, which had been most dear to you; that he had afterward been defended by you in a peril touching his life, aided in his canvass for the praetorship; that he had always hoped to have two closest friends, you by your kindness, me by his own. And if he could not prove this --- if that suspicion had so deeply lodged in you that it could in no way be torn out, if, in short, Italy were never to find rest from the levy, the city from arms, without Milo's destruction --- then he, doubtless, would not have hesitated to quit his country, born and bred as he is to such a course; but he would still, Magnus, call you to witness beforehand, as even now he does.

\ciceroSection{69} You see how varied and changeable is the course of life, how wandering and rolling is fortune, how great are the betrayals in friendships, how the pretences are fitted to the moment, how great in dangers are the flights of those closest to us, how great the cowardice. There will be --- there will surely be --- that time, and that day will one day dawn, when you, with your own fortunes safe, as I hope, but perhaps amid some upheaval of the common times --- and how often that happens we who have known it ought to know --- will long for the goodwill of your most faithful friend, the loyalty of a most steadfast man, and the greatness of spirit of the bravest man born since the human race began.

\ciceroSection{70} And yet, who would believe this --- that Cn. Pompey, most skilled in public law, in the custom of our forefathers, in the commonwealth itself, when the Senate had committed it to him to see that the commonwealth take no harm (and by that one little line the consuls have always been armed enough, even with no arms given them), would, with an army given him, with a levy given him, have awaited a trial in punishing the schemes of a man who by violence was abolishing the very courts? Judgement enough was passed by Pompey, enough, that those charges are falsely heaped upon Milo --- the man who carried the law by which, as I think, Milo ought to be acquitted by you, and, as all confess, may lawfully be.

\ciceroSection{71} But that he sits in that place, surrounded by those forces of the public guards, declares plainly enough that he is not bringing terror upon you --- for what would less befit him than to compel you to condemn a man on whom he could himself, both by the custom of our forefathers and by his own right, inflict punishment? --- but that he is a protection, that you may understand that, against yesterday's meeting, you are permitted to judge freely as you think.

\ciceroSection{72} Nor indeed, gentlemen, does the Clodian charge move me, nor am I so deranged, so ignorant and devoid of any sense of your feelings, as not to know what you think about the death of Clodius. And as to that, even if I were now unwilling to clear away the charge as I have cleared it, Milo would still be free, with impunity, to cry out openly and to boast with a glorious lie: ``I killed him --- I killed him: not Spurius Maelius, who, by lightening the corn-supply and at the cost of his own estate, fell under suspicion of aiming at kingship because he seemed to court the plebs too far, nor Tiberius Gracchus, who by sedition abrogated the magistracy of his colleague --- whose slayers filled the whole world with the glory of their name --- but the man'' (for he would dare to say it, since he had freed his country at his own peril) ``whose abominable adultery, in the most sacred shrines, the noblest women caught in the act;

\ciceroSection{73} the man whose punishment the Senate often decreed that the solemn rites must be expiated by; the man whom L. Lucullus, on oath, declared he had discovered, after holding examinations, to have committed wicked debauchery with his own full sister; the man who, by the arms of slaves, drove into exile the citizen whom the Senate, whom the Roman people, whom all nations had judged the preserver of the city and of the citizens' lives; the man who gave kingdoms, took them away, and parcelled out the whole world among whomever he pleased; the man who, after countless slaughters done in the Forum, drove a citizen of singular virtue and glory home by force and arms; the man to whom nothing was ever forbidden, neither in crime nor in lust; the man who burned the temple of the Nymphs in order to wipe out the public record of the census, stamped upon the public registers;

\ciceroSection{74} the man, in short, for whom by now there was no law, no civil right, no boundary-mark of property --- who sought other men's estates not by the chicanery of lawsuits, not by unjust claims and pledges, but by camps, by an army, by advancing standards; who tried to drive from their holdings, by arms and camps, not only the Etruscans --- for those he utterly despised --- but this P. Varius, a most brave and excellent citizen, our juror; who roamed the country houses and gardens of many men with architects and measuring-rods; who had bounded the hope of his holdings only by the Janiculum and the Alps; who, when he could not get a Roman knight of brilliant and bold character, M. Paconius, to sell him an island in the Prilian lake, suddenly hauled timber, lime, rubble, and sand in skiffs onto that island and, with the owner looking on from across the bank, did not hesitate to raise a building on land not his own; who said to this T. Furfanius --- what a man, immortal gods! ---

\ciceroSection{75} (for why should I speak of poor little Scantia, why of the young P. Aponius, both of whom he threatened with death unless they yielded him the possession of their gardens?) --- but who dared to say to T. Furfanius that, if he did not give him as much money as he demanded, he would carry a corpse into his house, by which scandal a man of such standing would have to be burned up; who threw his brother Appius --- a man bound to me by the most faithful goodwill --- out of the possession of an estate in his absence; who set about running a wall through his sister's forecourt, and laying foundations in such a way as not only to deprive his sister of her forecourt but of all entrance and threshold.

\ciceroSection{76} And yet even these things now seemed bearable, though he rushed alike upon the commonwealth, upon private men, upon those far off, upon those near, upon strangers, upon his own kin; but somehow, by long use, the city's incredible patience had grown hardened and callous. As for the things that were already at hand and impending --- how, pray, could you ever have warded them off, or borne them? Had he gained power --- I leave aside the allies, foreign nations, kings, tetrarchs; for you would have prayed that he might hurl himself upon them rather than upon your holdings, your homes, your money: money, do I say? From your children, so help me god, and from your wives he would never have held back his unbridled lusts. Do you think these things are invented --- things that lie open, that are known to all, that are established --- that he would have enrolled armies of slaves in the city, through whom he might seize the whole commonwealth and the private fortunes of all?

\ciceroSection{77} Therefore, if, holding a bloody sword, T. Annius were to cry out: ``Come, I beg you, and hear me, citizens! I have killed P. Clodius; his frenzies, which we could no longer curb by any laws, by any courts, I have driven back from your throats with this steel and this right hand, so that through me alone right and equity, the laws and liberty, decency and chastity might remain in the state'' --- would there really be cause to fear how the state would take it? For as it is, who is there who does not approve, who does not praise, who does not both say and feel that T. Annius alone, since the memory of mankind, has done the commonwealth the greatest service, and has filled the Roman people, all Italy, all nations with the highest joy? I cannot judge how great were those ancient joys of the Roman people; yet our age has by now seen many most illustrious victories of the greatest commanders, none of which brought a gladness either so lasting or so great.

\ciceroSection{78} Commit this to memory, gentlemen. I hope that you and your children will see many good things in the commonwealth; in each one of them you will always reckon thus: that, had P. Clodius been alive, you would have seen none of them. We have been led into the greatest hope, and, as I trust, the truest: that this very year, with this most eminent man as consul, men's licence checked, their cravings broken, the laws and the courts re-established, will be a year of safety for the state. Is there anyone, then, so deranged as to think that this could have come to pass while P. Clodius lived? And what of the things you hold as private and your own --- what right of perpetual possession could they have had, with a frenzied man as master? I am not afraid, gentlemen, that I may seem, inflamed by hatred born of my own enmities, to be venting these things upon him with more relish than truth. For although it ought to have been particular to me, yet he was so far the common enemy of all that my own hatred barely held its own amid the common hatred. It cannot be said enough, nor even conceived, how much of crime, how much of ruin was in that man.

\ciceroSection{79} Nay, attend in this way, gentlemen. Picture to your minds --- for our thoughts are free, and behold what they will in such a way that we discern it as we discern the things we see --- picture, then, in thought the image of this plight of mine: if I could bring it about that you acquit Milo, but only on this condition, that P. Clodius come back to life. Why have you taken fright in your faces? How would he, living, affect you, when, dead, he has struck you with an empty thought? What of this? If Cn. Pompey himself --- a man of such virtue and such fortune that he has always been able to do what no one but he could do --- if he, I say, had been able either to bring a bill of inquiry into the death of P. Clodius or to raise the man himself from the dead, which do you think he would rather have done? Even if for friendship's sake he should wish to call him back from the dead, for the commonwealth's sake he would not have done it. You sit, then, as avengers of the death of a man whose life, if you thought it could be restored through you, you would not wish restored; and a bill of inquiry has been carried touching the murder of one who, could he come to life again under that same law, would never have had that law carried at all. If, then, there were a slayer of this man, would he, in confessing it, fear punishment from those whom he had set free?

\ciceroSection{80} The Greeks award the honours of gods to those men who have slain tyrants --- the things I have seen at Athens, in the other cities of Greece! What rites of worship instituted for such men, what hymns, what songs! They are all but consecrated to immortality, both in religion and in remembrance --- and you, will you not only confer no honours upon the preserver of so great a people, the avenger of so great a crime, but even suffer him to be dragged off to execution? He would confess --- he would confess, I say, if he had done it, and would do so with a high heart and gladly: that he had acted for the freedom of all, a thing not only to be confessed, but truly to be proclaimed.

\ciceroSection{81} For if a man does not deny that for which he asks nothing but pardon, would he hesitate to avow that for which he ought to be claiming the rewards of praise? Unless, indeed, he supposes you find it more welcome that he was the defender of his own life than of yours --- especially when, by such a confession, were you willing to be grateful, he would attain the most ample honours. But if the deed found no favour with you --- though how could his own safety fail to find favour with anyone? --- yet still, if the courage of a most brave man fell out unwelcome to his fellow citizens, he would withdraw from an ungrateful state with a high and steadfast heart. For what could be more ungrateful than that all the rest should rejoice, and he alone should grieve, the very man on whose account the rest were rejoicing?

\ciceroSection{82} And yet this has always been the spirit in all of us, in crushing the betrayers of our country: that, since the glory would be ours, we should reckon the danger and the hatred ours as well. For what praise should be rendered to me myself --- when in my consulship I had dared so much for you and for your children --- if I had thought I could dare what I was attempting without the greatest struggles of my own? What woman would not dare to kill a wicked and ruinous citizen, if she did not fear the danger? But the man who, with hatred, death, and punishment set before him, defends the commonwealth no more slackly for that --- he is the man who must truly be reckoned a man. It is the mark of a grateful people to reward citizens who have deserved well of the commonwealth; it is the mark of a brave man not to be moved even by punishments to repent of having acted bravely.

\ciceroSection{83} Therefore let T. Annius use the same confession that Ahala used, that Nasica, that Opimius, that Marius, that we ourselves used; and, if the commonwealth were grateful, let him rejoice; if ungrateful, let him still, in a heavy lot, lean upon his own conscience. But the gratitude for this service, gentlemen, the Fortune of the Roman people and your own good fortune and the immortal gods reckon to be owed to themselves. Nor indeed can anyone judge otherwise, unless he holds that there is no power and no divine will --- a man whom neither the greatness of our empire moves, nor that sun, nor the motion of the sky and the constellations, nor the alternations and ordered courses of things, nor, what is greatest of all, the wisdom of our forefathers, who themselves cultivated the sacred rites, the ceremonies, the auspices with the utmost reverence, and handed them down to us, their descendants.

\ciceroSection{84} There is, there surely is that power; nor is it the case that, while in these bodies and in this frailty of ours there is something that has vigour and feeling, there is none in so vast and so glorious a movement of nature. Unless perhaps they think there is none because it does not show itself and is not seen --- just as if we could see our own mind, by which we are wise, by which we foresee, by which we do and say these very things, and plainly perceive of what kind it is, or where. That power itself, then, which has often brought to this city blessings and resources past belief, snuffed out and destroyed that pestilence: first casting into his mind that he should dare to provoke with violence and challenge with the sword a most brave man, and should be conquered by the man whom, had he conquered, he would have held impunity and licence everlasting.

\ciceroSection{85} It was not by human counsel, gentlemen, nor even by any ordinary one, but by the care of the immortal gods, that that deed was brought to pass. By Hercules, the very regions which saw that monster fall seem to have stirred themselves and to have kept their own right over him. For you now --- you hills and groves of Alba --- you, I say, I implore and call to witness; and you, ruined altars of the Albans, the partners and equals of the Roman people's sacred rites, which that man, headlong in his madness, had buried beneath the insane masses of his foundations, the most holy groves cut down and laid low: your sacred powers were then in their vigour, your might prevailed, the might which he had defiled with every crime. And you, holy Jupiter of Latium, from your high mountain --- you whose lakes and groves and bounds he had so often stained with every unspeakable defilement and crime --- you at last opened your eyes to punish him: to you, to you and in your sight, were those penalties paid, late, but just nonetheless, and owed.

\ciceroSection{86} Unless perhaps we shall say that this too came about by chance: that before the very shrine of the Good Goddess, which is on the estate of T. Sertius Gallus, a young man among the foremost in honour and distinction --- before the Good Goddess herself, I say --- when he had joined battle, he received that first wound by which he was to meet a most loathsome death; so that he should seem not acquitted by that abominable trial, but reserved for this signal punishment. Nor indeed did the same anger of the gods fail to cast this madness upon his followers: that without ancestral images, without dirge and games, without funeral rites, without lamentations, without eulogies, without a funeral, smeared with gore and mud, stripped of the solemnity of that last day which even enemies are wont to grant, he was burned, flung aside. It was not lawful, I believe, that the likenesses of the most illustrious men should bring any honour to that most loathsome parricide; nor that his death should be mangled in any place rather than in the one where his life had been condemned.

\ciceroSection{87} Harsh, so help me god, and cruel did the Fortune of the Roman people seem to me by now, which had suffered him for so many years to trample upon this commonwealth. He had defiled the most holy rites with debauchery, he had shattered the gravest decrees of the Senate, he had openly bought himself off from the jurors with money, he had harried the Senate in his tribunate, he had undone what had been done by the consensus of all the orders for the safety of the commonwealth, he had driven me from my country, he had plundered my goods, he had burned my house, he had harried my children and my wife, he had declared an unspeakable war upon Cn. Pompey, he had wrought the slaughter of magistrates and of private men, he had burned my brother's house, he had laid Etruria waste, he had cast many men out of their homes and their fortunes; he pressed on, he bore down; the state, Italy, the provinces, the kingdoms could not contain his madness; already laws were being cut into bronze in his house that would make us over to our own slaves; there was nothing belonging to anyone --- nothing, at least, on which he had set his heart --- that he did not reckon would be his this year.

\ciceroSection{88} No one stood in the way of his designs but Milo. The one man who could stand in his way he reckoned bound to himself by a new return to favour; he kept saying that Caesar's power was his own; he had despised the spirit of good men in my own downfall: Milo alone bore down upon him. At this point the immortal gods, as I said above, put it into the mind of that desperate madman to lay an ambush for Milo. There was no other way for that pest to perish; the commonwealth would never have avenged itself upon him by its own right. The Senate, I suppose, would have curbed him as praetor. Not even when it used to do so had it made any headway against him as a private man.

\ciceroSection{89} Or would the consuls have been brave in keeping a praetor in check? First, once Milo was killed, he would have had his own consuls; next, what consul would be brave in dealing with that praetor, through whom he remembered that the consular authority of a tribune had been most cruelly harried? He would have crushed everything, possessed everything, held everything; by a new law --- one found among his papers along with the rest of the Clodian laws --- he would have made our slaves into his own freedmen; in short, had not the immortal gods driven him into the mind to attempt, an effeminate man, the killing of a most brave one, you would have had no commonwealth today.

\ciceroSection{90} Would that man as praetor, that man indeed as consul --- if only these temples and these very walls could have stood so long while he lived, and could have awaited his consulship --- would he, in fine, while alive, have done no harm, when, once dead, one of his own followers set fire to the Senate-house on his account? And what have we seen more wretched, more bitter, more grievous than this? That a temple of holiness, of greatness, of intelligence, of public counsel, the head of the city, the altar of our allies, the haven of all nations, the seat granted by the whole people to a single order, should be set ablaze, demolished, defiled --- and that this should be done not by an ignorant mob, wretched as that itself would be, but by a single man? When he who burned the corpse dared so much for the dead man, what would the standard-bearer not have dared for the living? Into the Senate-house above all places he flung him, so that dead he should burn down what alive he had overthrown.

\ciceroSection{91} And there are those who complain of the Appian Way, yet say nothing of the Senate-house; who suppose that the Forum could have been defended against him while he breathed, when his corpse the Senate-house could not withstand. Raise him up, raise him up, if you can, from the dead: will you break the onset of the living, you who scarcely endure the fury of his unburied body? Or did you withstand those men who ran together with firebrands to the Senate-house, with axes to the temple of Castor, who flitted through the whole Forum with swords? You saw the Roman people cut down, the assembly broken up by swords, while M. Caelius, tribune of the plebs, was being heard in silence --- a man both in public life most brave and in the cause he had taken up most firm, devoted to the will of good men and the authority of the Senate, and, in this matter of Milo's --- call it his unpopularity or his fortune --- of a loyalty singular, godlike, past all belief.

\ciceroSection{92} But now I have said enough about the case, and outside the case perhaps even too much. What remains but that I beg and beseech you, gentlemen, to grant to a most brave man the pity which he himself does not implore, but which I, against his will, both implore and demand? Do not, if amid the weeping of us all you have caught sight of no tear from Milo, if you see his face always the same, his voice, his speech steady and unaltered --- do not for that reason spare him the less: I rather think he should perhaps be helped all the more. For if, in gladiatorial combats and in the lot and condition of men of the lowest sort, we are wont even to hate the timid and the suppliant and those who beg to be allowed to live, while we long that the brave and spirited, who fiercely offer themselves to death, be spared --- and if we pity more those who do not seek our pity than those who clamour for it --- how much more ought we to do this in the case of most brave citizens!

\ciceroSection{93} For my part, gentlemen, I am unmanned and undone by these words of Milo's, which I hear without ceasing and am present to hear every day. ``Farewell,'' he says, ``farewell, my fellow citizens; may they be safe, may they flourish, may they prosper; may this glorious city stand, my country dearest to me, however it shall have deserved of me; may my fellow citizens, since I may not enjoy it with them, enjoy the commonwealth in peace without me, yet on my account. I will yield and go. If I may not enjoy a good commonwealth, at least I shall be free of a bad one; and the first well-ordered and free state I reach, in that I will find my rest.''

\ciceroSection{94} ``O labours of mine,'' he says, ``undertaken in vain! O deceiving hopes! O empty thoughts of mine! When, as tribune of the plebs, with the commonwealth crushed, I gave myself to the Senate, which I had received as good as extinguished, to the Roman knights whose strength was feeble, to the good men who had thrown away all authority before the arms of Clodius --- could I ever have thought that the protection of good men would fail me? When I had given you back --- for he speaks with me very often --- to your country, could I have thought there would be no place for me in my own? Where now is the Senate which we followed? Where the Roman knights --- those, he says, those of yours? Where the zeal of the towns? Where the voices of Italy? Where, in fine, your own voice and defence, M. Tullius, which has been a help to so many? To me alone, who have so often offered myself to death for you, can it bring no aid?''

\ciceroSection{95} And he says these things, gentlemen, not, as I do now, weeping, but with the very face you see. For he denies --- he denies that he did what he did for ungrateful citizens; for timid ones, who eye every danger about them, he does not deny it. The common people and the lowest crowd, which under Clodius's leadership was threatening your fortunes --- that he made safer, that your own lives might be the more secure; so that he not only swayed it by his courage but soothed it with his three patrimonies; nor does he fear that, having appeased the people with shows, he has failed to win you over by his singular services to the commonwealth. The Senate's goodwill toward him has often been seen clearly in these very times; but your own kindnesses, and those of the men of your order --- your meetings, your zeal, your talk --- whatever course Fortune shall take, he says he will carry away with him.

\ciceroSection{96} He remembers too that the herald's voice alone was wanting to him --- which he wished for least of all --- but that by all the votes of the people, the one thing he had desired, he was declared consul; and now at last, if these arms are to be turned against him, it is the suspicion of a deed, not the charge of one, that stands against him. He adds this, which is surely true: that brave and wise men are wont to pursue not so much the rewards of right deeds as the right deeds themselves; that he has done nothing in his life but what was most glorious, since there is nothing more excellent for a man than to free his country from peril.

\ciceroSection{97} Happy, he says, are those for whom this deed has been an honour at the hands of their fellow citizens; yet not even those are wretched who have surpassed their fellow citizens in a benefit conferred. But still, of all the rewards of virtue --- if account were to be taken of rewards --- the most ample reward is glory; this alone it is that consoles the brevity of life with the remembrance of posterity, that brings it about that, though absent, we are present, and, though dead, we live; this, in fine, it is by whose steps men seem even to ascend into heaven.

\ciceroSection{98} ``Of me,'' he says, ``the Roman people will always speak, always all the nations; no length of time will ever fall silent. Nay, at this very moment, when all the torches of hatred against me are being thrust beneath me by my enemies, still in every gathering of men we are honoured with the giving of thanks and the offering of congratulations and in all men's talk. I pass over the festal days of Etruria, both held and instituted. This is the hundredth day, I think, and one more, since the death of P. Clodius; and to wherever the bounds of the Roman people's empire reach, not only the report of that deed by now, but the rejoicing in it, has travelled abroad. And so where this body of mine may be,'' he says, ``I do not trouble myself, since in all lands the glory of my name already moves and will forever dwell here.''

\ciceroSection{99} These things you have often said to me with these men absent; but these things I say to you, Milo, with these same men listening: you indeed, since you are of such a spirit, I cannot praise enough; but the more godlike that virtue of yours, the greater the grief with which I am torn from you. Nor indeed, if you are wrenched from me, is even that one complaint left me for consolation --- that I might be angry with those from whom I had received so great a wound. For it is not my enemies who will wrench you from me, but my dearest friends; not men who have ever deserved ill of me, but men who have always deserved supremely well. You will never, gentlemen, brand me with a grief so great --- and yet what grief could be so great? --- no, not even this very one, as to make me forget how highly you have always valued me. But if that forgetting has seized you, or if you have taken some offence in me, why is it not rather paid for by my head than by Milo's? For I shall have lived gloriously, if some accident befalls me before I see this great evil come to pass.

\ciceroSection{100} As it is, one consolation upholds me: that toward you, T. Annius, no office of love, of zeal, of devotion has been wanting on my part. For you I sought the enmities of the powerful; for you I have often flung my own body and life against the weapons of your enemies; for you I have cast myself down as a suppliant before very many; my goods, my fortunes and those of my children I have brought into partnership with your trials; this very day, in fine, if any force is in readiness, if any struggle for life is at hand, I claim it for my own. What now remains? What have I left to do for your services to me but to count whatever fortune shall be yours as mine? I do not refuse it, I do not shrink from it; and I beseech you, gentlemen, either to crown the kindnesses you have heaped upon me by his preservation, or to see that they are doomed to fall with this same man's destruction.

\ciceroSection{101} By these tears Milo is not moved --- he is of a certain incredible strength of mind --- he thinks exile to be there, where there is no place for virtue, and death to be the end of nature, not a punishment. Let him be of the spirit he was born with: but what of you, gentlemen? With what spirit, in the end, will you be? Will you keep the memory of Milo, and cast out the man himself? And will there be any place in all the earth worthier to receive this virtue than the one that brought it forth? You --- you I appeal to, bravest of men, who have shed much blood for the commonwealth; you, I say, in the peril of an unconquered citizen, I appeal to, centurions, and you, soldiers: while you not only look on but stand armed and guard this court, shall this virtue, so great, be driven from this city, banished, cast out?

\ciceroSection{102} O wretched that I am! O unhappy that I am! You, Milo, could call me back to my country through these men; shall I not be able, through these same men, to keep you in your country? What shall I answer my children, who think you their second parent? What shall I answer you, my brother Quintus, who are now away, the partner with me of those times? That I could not protect Milo's safety through the same men through whom he had preserved ours? And in what cause could I not? One that is welcome to all the nations. From what men could I not? From those who took most comfort in the death of P. Clodius. With whom pleading? Me.

\ciceroSection{103} What crime so great did I conceive, gentlemen, or what monstrous deed did I commit against myself, when I tracked out the proofs of our common ruin, laid them bare, brought them forth, stamped them out? All my griefs, and my family's, well up from that one source. Why did you wish me to return? Was it that, with me looking on, the men by whom I was restored should be driven out? Do not, I beg you, suffer my return to be more bitter to me than that very departure was. For how can I think myself restored, if I am torn away from the men through whom I was restored? Would that the immortal gods had brought it about --- by your leave, my country, let me say it, for I fear lest what I shall say dutifully for Milo I shall say criminally against you --- would that P. Clodius were not only alive, but even praetor, consul, dictator, rather than that I should look upon this spectacle!

\ciceroSection{104} O immortal gods! A brave man, gentlemen, and one to be preserved by you! ``By no means, by no means,'' he says; ``rather let him pay the penalties he owed; let us undergo, if it must be so, penalties not owed.'' Shall this man, born for his country, die anywhere but in his country, or, perhaps, for his country? Will you keep the memorials of his spirit, and suffer his body to have no tomb in Italy? Will anyone, by his own vote, drive from this city the man whom, once you have driven him out, all cities will call to themselves?

\ciceroSection{105} O happy that land which shall receive this man! Ungrateful, if it cast him out! Wretched, if it lose him! But let there be an end; for I can no longer speak for tears, and this man forbids that he be defended with tears. I beg and beseech you, gentlemen, that in casting your votes you dare to vote what you truly judge. Your courage, your justice, your good faith --- believe me --- he will most of all approve, who in choosing the jurors picked out each best and wisest and bravest man.
